\documentclass[11pt]{article}

% --------------------
% Packages
% --------------------
\usepackage[a4paper,margin=1in]{geometry}
\usepackage{amsmath,amssymb,amsfonts}
\usepackage{hyperref}

% --------------------
% Document
% --------------------
\begin{document}

\title{Time Complexity Analysis of APEX and LFTJ}
\author{}
\date{}
\maketitle

\paragraph{\textbf{Background.}}
A basic graph pattern (BGP) query can be reduced to a natural join over a set of relations derived from triple patterns.
Each triple pattern $(s,p,o)$ is mapped to a relation by fixing all constant positions and treating the remaining variable positions as attributes.
For example, when the predicate $p$ is fixed and both $s$ and $o$ are variables, the triple pattern corresponds to a binary relation $R(s,o)$.
If either the subject or the object is a constant, the resulting relation becomes unary.

Under this relational interpretation, evaluating a BGP query amounts to computing a natural join over the relations induced by all triple patterns.
Shared variables across triple patterns correspond to join attributes.

Consider a query $Q(v_0,\dots,v_{k-1})$ with $k$ variables and a fixed variable ordering $\Phi = [v_0,\dots,v_{k-1}]$. 
Let $N_{\max}$ denote the maximum cardinality of input relations.

Under ordering $\Phi$, the Leapfrog Triejoin (LFTJ) algorithm runs in
$$
T_{\text{LFTJ}} = O\!\left(\left(\sum_{i=0}^{k-1} C_i\right)\log N_{\max}\right).
$$
Here, $N_{\max}$ is the maximum size of an input relation and $C_i$ denotes the \emph{sum-min count} for variable $v_i$. 
The sum-min count is defined as the sum over all prefix bindings of $v_i$, where a prefix binding corresponds to a binding combination generated for the preceding variables $[v_0, \ldots, v_{i-1}]$ during LFTJ execution.
For each such prefix binding, $C_i$ counts the minimum list length among all lists participating in the intersection for $v_i$.

Since the size of an intersection result is bounded by the length of the shortest participating list, $C_i$ upper-bounds the total number of candidate values generated for variable $v_i$.

\paragraph{\textbf{Intersection Time Complexity.}}
We analyze that, under the same variable order $\Phi=[v_0, \ldots, v_{k-1}]$, LFTJ and APEX perform the same intersection operations when dependency-source detection and Bloom filters are not applied.
% {\color{red} [check ``the same'' or maybe less than? Whether APEX computes the intersection for $v_i$ at most once, then reuse?].}
% 如果不使用之后的两个优化技巧，是一样的。获取一个变量的交集结果后，会reuse

\paragraph{Initial Case ($i = 0$).}
For the first variable $v_0$, LFTJ considers only the empty prefix binding. Similarly, in APEX, the CPG is empty at initialization and contains no nodes or candidates of variables; Algorithm~2 starts from the empty prefix binding $\sigma = \emptyset$ and invokes the \texttt{Traverse} procedure.

At this stage, both methods perform exactly one intersection operation. They retrieve the value lists for $v_0$ from all triple patterns that contain $v_0$, subject to the existing constant constraints, and compute their intersection. Because the input lists to the intersection are identical, LFTJ and APEX produce the same intersection result for $v_0$.

\paragraph{Inductive Step.}
Assume that for all variables preceding $v_i$, LFTJ and APEX perform the same number of intersection operations and the results of intersections is also the same. We now consider the $i$-th variable $v_i$.

In LFTJ, the algorithm employs DFS following the variable order $\Phi$ to enumerate all valid bindings of all variables.
Due to the DFS traversal, the number of binding combinations processed for $v_0, \ldots, v_{i-1}$ equals the sum of the sizes of the intersection results obtained for $v_{i-1}$.

In APEX, Algorithm~2 likewise uses DFS to traverse the intersection results stored in the CPG, thereby enumerating all possible bindings of variables $v_0, \ldots, v_{i-1}$. Thus, the number of binding combinations is again equal to the sum of the sizes of the intersection results of $v_{i-1}$.

Furthermore, since APEX and LFTJ share the same variable ordering and traverse same intersection results of preceding variables, the binding combinations generated by the two methods are identical in both cardinality and content.
For each such binding combination, both LFTJ and APEX perform exactly one intersection operation for variable $v_i$. Therefore, at the $i$-th variable, APEX and LFTJ execute the same intersection operations.

\paragraph{Conclusion.}
By induction, it follows that under the given variable ordering and assumptions, APEX and LFTJ perform exactly the same intersection operations at each variable throughout the entire execution.
Hence, the total time complexity of intersection operations in APEX is the same as LFTJ, and can be expressed as:
$$
T_{\text{join}}(\text{APEX})
= O\!\left(\Big(\sum_{i=0}^{k-1} C_i\Big)\log N_{\max}\right),
$$

\paragraph{\textbf{Additional Overhead of APEX.}}
Compared with LFTJ using the same variable order, the additional overhead of APEX mainly comes from two types of operations: (i) traversing the CPG to enumerate prefix bindings; and (ii) materializing the candidate values corresponding to each prefix binding and storing it into the CPG.

As shown above, APEX and LFTJ with the same variable order perform identical intersection operations and produce intersection results of the same total size at each variable. 
The total size of the intersection results produced by LFTJ at the $i$-th variable is at most $C_i$. 
Since the CPG stores exactly the candidate values obtained from these intersection operations for all variables, the cost of traversing the CPG is $O(\sum_{i=0}^{k-1} C_i).$
As APEX processes the $k$ variables one by one, such a traversal may occur at most $k$ times. Thus, the total traversal cost is
$$
T_{\text{traverse}}
= O\!\left(k\sum_{i=0}^{k-1} C_i\right).
$$

For each prefix binding $\sigma$ of variable $v_i$, APEX uses $\sigma$ as the key and stores the intersection result as the value in the CPG. 
The number of prefix bindings for $v_i$ equals the total size of the intersection results produced for the preceding variable, and is therefore bounded by $C_{i-1}$. When $i=0$, there is only one intersection operation, corresponding to a single key. Hence, the total number of insertions is bounded by $1 + \sum_{i=1}^{k-1} C_{i-1}.$
Accordingly, the storage cost is
$$
T_{\text{store}}
= O\!\left(1 + \sum_{i=1}^{k-1} C_{i-1}\right)
= O\!\left(\sum_{i=1}^{k-1} C_{i-1}\right).
$$

\paragraph{\textbf{Overall Time Complexity of BFS Execution for LFTJ.}}
Since the number of variables $k$ can be treated as a constant, the overall time complexity of APEX is given by
$$
\begin{aligned}
T_{\text{APEX}}
&= T_{\text{join}} + T_{\text{traverse}} + T_{\text{store}} \\
&= O\!\left(\Big(\sum_{i=0}^{k-1} C_i\Big)\log N_{\max}\right)
\;+\; O\!\left(k\sum_{i=0}^{k-1} C_i\right)
\;+\; O\!\left(\sum_{i=1}^{k-1} C_i\right) \\
&= O\!\left(\Big(\sum_{i=0}^{k-1} C_i\Big)\log N_{\max}\right).
\end{aligned}
$$
Based on the formal time complexity analysis, we show that BFS execution of LFTJ using a CPG incurs additional overhead compared to standard LFTJ, but these overheads are bounded and do not affect the worst-case asymptotic complexity. 

\end{document}
